\documentclass[a4paper]{article}
\usepackage[pdftex]{graphicx}
\usepackage[utf8]{inputenc}
\usepackage{enumerate}
\usepackage{icomma}
\usepackage{amssymb}
\usepackage{tikz}
\usepackage{href-ul}
\hypersetup{
	colorlinks=true,
	linkcolor=blue,
	urlcolor=blue}
\usepackage{geometry}
\geometry{a4paper, top=15mm, left=15mm, right=15mm, bottom=15mm,
	headsep=10mm, footskip=12mm}

% Start the document
\begin{document}
\begin{tikzpicture}(10,3)
	\draw[ultra thick](2,0) --(10,0) -- (10,1.3) --(2,1.3) -- (2,0);
	\draw[fill=black](2,0)-- (0,.6) -- (2,.6) -- (2,0);
	\node at (6,.6) {\large Lösungsblatt von \href{https://www.okuyakl.de}{www.okuyakl.de}};
\end{tikzpicture}
\vspace{0.5 cm}

\noindent{\bf Aufgabe 1.}\\
\renewcommand{\arraystretch}{2.7}
\begin{tabular}{|p{40 pt}|p{60 pt}|p{40 pt}|p{40 pt}|p{70 pt}|p{70 pt}|p{100 pt}|}
\hline
 & Funktionstyp & Potenz $n$ &  Faktor $a$ & Definitionsmenge $\mathbb{D}=$ & Wertemenge $\mathbb{W}=$ &  Funktionsterm  \\
\hline
$f_1(x)$ & Gerade & 1 & -0,25 & $\mathbb{R}$ & $\mathbb{R}$ & $f(x)=-0,25x$\\
\hline
$f_2(x)$ & Hyperbel & $-1$ & 4 & $\mathbb{R}\setminus \{0\}$ & $\mathbb{R}\setminus \{0\}$ & $f(x)=4x^{-1}={4 \over x}$\\
\hline
$f_3(x)$ & Parabel & 2 & 3 & $\mathbb{R}$ & $[0;\infty[$ & $f(x)=3x^2$\\
\hline
$f_4(x)$ & Wurzel\-funktion & ${1\over 2}$ & 1 & $[0;\infty[$ & $[0;\infty[$ & $f(x)=x^{1\over 2}=\sqrt{x}$\\
\hline
$f_5(x)$ & Wende\-parabel& 3 & $-{1\over 125}$ & $\mathbb{R}$  &$\mathbb{R}$ & $f(x)=-{1\over 125}x^3$\\
\hline
$f_6(x)$ & Hyperbel & $-2$ & $-1$ &  $\mathbb{R}\setminus \{0\}$ & $]-\infty;0[$ &  $f(x)=- x^{-2} = -{1\over x^2} $\\
\hline
\end{tabular}

\begin{center}
	\includegraphics[width=7 cm]{../../viecher/endcomic.pdf}
	
	Hier geht es zurück zum \href{https://www.okuyakl.de/math/m10potiL021/aa021.pdf}{Aufgabenblatt}
\end{center}

%Ende
\end{document}